% !TEX root = ../Thesis.tex
\chapter{INTRODUCTION}
\section{Background and motivation}
Urban development changes how rainfall moves across the landscape. Roads, roofs, and parking areas reduce infiltration and increase the amount of rainwater that must be handled by drainage systems. Managing stormwater is therefore an important element of flood resilience and sustainable urban planning \parencite{epaGreenInfrastructure}.

\section{Problem statement}
Technical stormwater information can be difficult for non-specialist learners to interpret. An online learning tool can make the topic more approachable through brief explanations, interactive questions, and immediate feedback. The educational value of such a tool must be evaluated rather than assumed.

\section{Objectives}
The objective of this project is to further develop and test an accessible and engaging web-based learning experience addressing stormwater and urban flood resilience. The work includes development of learning content and evaluation of usability, accessibility, and understanding.

\section{Scope and report structure}
This version focuses on an initial browser-based prototype, three introductory learning modules, and a documented evaluation plan. It does not implement a hydrological simulation or municipal decision-support system. Chapters 2--4 present the background, methods, and implementation; Chapters 5--7 cover evaluation, discussion, and conclusions.
